\documentclass[12pt]{article}

% Layout
\usepackage[margin=2.5cm]{geometry}
\linespread{1.2}

% Math
\usepackage{amsmath}

% Figures and tables
\usepackage{graphicx}
\usepackage{float}
\usepackage{caption}
\usepackage{booktabs}
\usepackage{subcaption}
\usepackage{comment}
\usepackage{algorithm}
\usepackage{algpseudocode}

% Better URL breaking
\usepackage{xurl}

% Bibliography
\usepackage[backend=biber,style=numeric,backref=true]{biblatex}
\addbibresource{references.bib}

% Hyperlinks
\usepackage[colorlinks=true,citecolor=black,linkcolor=black,urlcolor=black]{hyperref}

\title{Advanced Machine Learning}
\author{
Project 1 Report\\[1em]
Wojciech Grunwald, Piotr Jurczyk and Krzysztof Krawiec\\
Warsaw University of Technology
}
\date{April 2026}

\begin{document}

\maketitle


\begin{abstract}
\noindent The aim of the project is to analyze a logistic regression model in a situation where the training 
dataset contains observations with missing labels.
\end{abstract}

\newpage 

{\footnotesize
\tableofcontents
}
\clearpage

\section{Introduction}

In this project, we investigate the performance of logistic regression in the presence of missing labels. We consider different missing data mechanisms (MCAR, MAR, MNAR) and analyze their impact on model training and prediction.

To address this problem, we implement label-masking procedures, an optimization algorithm based on FISTA, and an extended logistic regression model capable of handling partially observed data. The proposed methods are evaluated on multiple datasets with varying characteristics to assess their robustness and generalization ability.

\section{Preparing Datasets}

We selected four datasets representing different classification problems to evaluate model robustness and generalization.

\begin{enumerate}
    \item \href{https://archive.ics.uci.edu/dataset/19/car+evaluation}{\textbf{Car Evaluation Dataset}} consists of 1728 instances described by 15 features. The task is to classify whether a car is acceptable or not based on attributes such as \textit{number of doors}, \textit{luggage boot size}, and \textit{safety level}.
    
    \item \href{https://www.openml.org/search?type=data&sort=runs&status=active&qualities.NumberOfClasses=%3D_2&qualities.NumberOfFeatures=between_100_1000&id=1485}{\textbf{Madelon Dataset}} is an artificial dataset created for the NIPS 2003 feature selection challenge. It is particularly challenging due to its highly multivariate and non-linear structure. The dataset comprises 2600 instances with 500 features, many of which are redundant or irrelevant.

    \item \href{https://archive.ics.uci.edu/dataset/320/student+performance}{\textbf{Student Performance Dataset}} aims to predict whether a student passes a mathematics course. It includes academic results along with demographic, social, and school-related attributes. The dataset contains 395 instances described by 26 features.

    \item \href{https://www.openml.org/search?type=data&sort=runs&status=active&qualities.NumberOfClasses=%3D_2&qualities.NumberOfFeatures=between_100_1000&id=40536}{\textbf{Speed Dating Dataset}} addresses a binary classification problem of predicting whether a match occurs between participants. The data was collected during experimental speed dating events and includes features describing participants’ preferences, attributes, and interactions. It consists of 8378 instances and 329 features.
\end{enumerate}


\section{Implementing Missing Data Generators}

We implemented four mechanisms for masking the target vector $Y$ within a \texttt{DataGenerator} class.

\subsection{MCAR}

In MCAR (Missing Completely At Random), missingness is independent of both $X$ and $Y$:
\begin{align*}
    P(S = 1 \mid X, Y) = c,
\end{align*}
where $c \in [0,1]$. Missing values are generated by sampling a binary indicator vector $S$ from a Bernoulli distribution with parameter $c$. The observed target $y_{\text{obs}}$ is obtained by replacing entries corresponding to $S=1$ with a placeholder value (e.g., $-1$).

\subsection{MAR1}

In MAR1 (Missing At Random), missingness depends on a single feature $X_j$, where $j \sim \mathcal{U}(0, p-1)$:
\begin{align*}
    P(S = 1 \mid X, Y) = P(S = 1 \mid X_j).
\end{align*}
A feature index is randomly selected, and its values are transformed using a logistic function to obtain masking probabilities. These probabilities are then used to sample the indicator vector $S$, which determines which labels are replaced by a placeholder.

\subsection{MAR2}

In MAR2, missingness depends on all features:
\begin{align*}
    P(S = 1 \mid X, Y) = P(S = 1 \mid X).
\end{align*}
A score is computed as a normalized sum of all feature values and transformed via a logistic function to obtain masking probabilities. The indicator vector $S$ is then sampled accordingly, and missing labels are introduced in the same way as before.

\subsection{MNAR}

In MNAR (Missing Not At Random), missingness depends on both $X$ and $Y$:
\begin{align*}
    P(S = 1 \mid X, Y).
\end{align*}
A combined score based on feature values and the target variable is computed and passed through a logistic function to obtain masking probabilities. The indicator vector $S$ is sampled from these probabilities, and labels corresponding to $S=1$ are replaced with a placeholder.



\newpage
\section{FISTA Algorithm}


\subsection{Overview and Motivation}
The Fast Iterative Shrinkage-Thresholding Algorithm (FISTA) is a first-order optimization method designed specifically for composite convex optimization problems. These problems typically involve minimizing an objective function that can be decomposed into two distinct components: a smooth, continuously differentiable loss function and a possibly non-smooth, convex penalty term:
$$
    \min_{\mathbf{w}, b} F(\mathbf{w}, b) = f(\mathbf{w}, b) + g(\mathbf{w})
$$
For standard binary classification, we assume the input labels are $y^{(i)} \in \{0, 1\}$. To formulate the logistic loss using margin-based notation, we map these to signed labels $\tilde{y}^{(i)} \in \{-1, 1\}$. Given feature vectors $\mathbf{x}^{(i)}$, the objective components are defined as:
$$
    f(\mathbf{w}, b) = \frac{1}{n} \sum_{i=1}^n \log\left(1 + \exp\left(-y^{(i)}(\mathbf{w}^T \mathbf{x}^{(i)} + b)\right)\right), \quad \quad g(\mathbf{w}) = \lambda \|\mathbf{w}\|_1
$$
where $\lambda > 0$ acts as the regularization hyperparameter. To iteratively minimize the objective function formulated above, we employ the standard FISTA procedure outlined in Algorithm~1.
\begin{algorithm}[H]
\caption{FISTA with constant stepsize (Beck \& Teboulle, 2009)}
\begin{algorithmic}[1]
\Require $L = L(f)$ \Comment{A Lipschitz constant of $\nabla f$}
\Require $\mathbf{x}_0 \in \mathbb{R}^n$
\State \textbf{Step 0:} Initialize $\mathbf{y}_1 = \mathbf{x}_0$, $t_1 = 1$.
\For{$k = 1, 2, \dots$}
    \State $\mathbf{x}_k = p_L(\mathbf{y}_k)$
    \State $t_{k+1} = \frac{1 + \sqrt{1 + 4t_k^2}}{2}$
    \State $\mathbf{y}_{k+1} = \mathbf{x}_k + \left(\frac{t_k - 1}{t_{k+1}}\right)(\mathbf{x}_k - \mathbf{x}_{k-1})$
\EndFor
\end{algorithmic}
\end{algorithm}

\subsection{Implementation Details}
In our programmatic implementation, the iterative optimization relies strictly on vectorized mathematical operations. To align with FISTA notation, we denote the model weights as $\mathbf{x}$ and the intercept as $b$. We initialize the weights $\mathbf{x}_0$ and intercept $b_0$ to zero. To facilitate acceleration, auxiliary momentum variables $\mathbf{y}_1 = \mathbf{x}_0$, a bias momentum $y_{1,b} = b_0$, and a time step $t_1 = 1$ are introduced. The implementation then strictly follows the algorithm. Let us focus on how $\mathbf{x}_k$ is calculated in each iteration, using the soft-thresholding operator:
    $$
        \mathbf{x}_k = p_L(\mathbf{y}_k)=\mathcal{T}_{\alpha \lambda}(\mathbf{x}_{\text{temp}})=\mathcal{T}_{\alpha \lambda}(\mathbf{x}_{\mathbf{y}_{k}} - \alpha \nabla_{\mathbf{y}_k} f))
    $$
    where the proximal operator is defined as:
        $ \mathcal{T}_{\gamma}(z) = \text{sign}(z) \cdot \max(|z| - \gamma, 0)$,     
    which leads to the element-wise update:
        $$ \mathbf{x}_k = \text{sign}(\mathbf{x}_{\text{temp}}) \odot \max(|\mathbf{x}_{\text{temp}}| - \alpha \lambda, 0)     $$
    Because the intercept $b$ remains unregularized, it is updated purely via the gradient step: $$b_k = y_{k,b} - \alpha \nabla_{y_{k,b}} f.$$
Note, that to determine the sensitivity of the loss $f_i$ with respect to the momentum weights $\mathbf{y}_k$, we first define the logit as $z_i = (\mathbf{x}^{(i)})^T \mathbf{y}_k + y_{k,b}$. Consequently:
        $ \frac{\partial z_i}{\partial \mathbf{y}_k} = \mathbf{x}^{(i)}     $.
    Applying the chain rule to the logistic loss function $f_i$ with respect to the logit $z_i$ yields:
    $$
        \frac{\partial f_i}{\partial z_i} = \frac{1}{1 + \exp\left(-y^{(i)} z_i\right)} \cdot \exp\left(-y^{(i)} z_i\right) \cdot \left(-y^{(i)}\right).
    $$
    which can be simplified to the residual error term $\delta_i$:
    $$
        \delta_i = \frac{\partial f_i}{\partial z_i} = \left(\sigma\left(y^{(i)} z_i\right) - 1\right) y^{(i)}
    $$
Using the chain rule, we obtain the gradient for a single observation as
        $ \frac{\partial f_i}{\partial \mathbf{y}_k} = \delta_i \cdot \mathbf{x}^{(i)}     $.
To compute the total gradient across all $n$ observations, we average the individual gradients. The iterative process continues until a predefined convergence criterion is satisfied, specifically when the Euclidean norm of the parameter update, $\|\mathbf{x}_k - \mathbf{x}_{k-1}\|_2$, falls beneath a strict tolerance threshold.


\subsection{Hyperparameter Tuning}
To pick the best regularization coefficient $\lambda$, our implementation employs a grid search strategy combined with hold-out validation. The algorithm operates over a logarithmically spaced sequence of candidate values and computes the corresponding model weights. Then it is evaluated on a dedicated, hold-out validation set using a user-specified performance metric. The $\lambda$ that maximizes the validation score is selected as the optimal hyperparameter, and its associated coefficients are retained as the final model parameters.

\subsection{Prediction, Classification and Model Evaluation}
Following optimization, the framework provides standard inference and evaluation methods:
\begin{itemize}
    \item \texttt{predict\_proba} method computes the posterior class probabilities
    \item \texttt{predict} method converts the continuous probability estimates into binary class labels
    \item \texttt{plot}: Visualizes the selected performance metric as a function of $\lambda$
    \item \texttt{plot\_coefficients}): Illustrates the trajectory of individual feature coefficients as $\lambda$ varies.
\end{itemize}




\section{UnlabeledLogReg}


\end{document}
